\section{From curves to claims}
\label{sec:gates}

A rising curve is not a phase transition. This section states the decision
procedure that converts condition-level estimates into one of a small set of
named outcomes, including outcomes that block the strong reading.

\subsection{Unit of independence and uncertainty}

A \emph{condition} is one point of a family's design grid: a control value, a
size, and any registered secondary knobs. A \emph{run} is one condition under
one outer seed, and the run is the unit of independence. Inner stochastic
evaluations are averaged within a run before uncertainty is computed, so no
inner replicate is ever promoted to an outer sample. Intervals are two-sided
$95\%$ Student-$t$ over the realised outer-seed count, which is carried through
aggregation and printed on every figure.

\subsection{Three finite-size signatures}

For each family we compute three independent signatures from the same
condition-level data.

\paragraph{Sharpening.} At each size we locate the interior maximum of
$\chi_N$. A maximum at either end of the sampled control range is
\emph{censored} rather than reported as a critical point: the grid does not
contain the boundary, and naming the endpoint would manufacture one. Across the
sizes that do have an interior peak we fit
\begin{equation}
  \chi_N^{\rm peak}\propto N^{\,\omega},
  \label{eq:sharpening}
\end{equation}
and $\omega>0$ is the signature that the response diverges rather than
saturating.

\paragraph{Location.} The peak position itself drifts with size as
\begin{equation}
  c_{\rm peak}(N)=c_\infty+A\,N^{-1/\nu},
  \label{eq:drift}
\end{equation}
so the boundary of the infinite system is obtained by extrapolation, never by
reading off the largest size available. A fit is admissible only if $c_\infty$
is finite, precisely estimated, and interior to the sampled range. Peaks that
march monotonically out of the sampled window produce a formally convergent
least-squares fit with an enormous intercept and an enormous standard error;
that situation is recorded as an \emph{unresolved boundary}, which is a
different --- and weaker --- statement than a measured one.

\paragraph{Collapse.} Independently, we scan $(c^\ast,1/\nu)$ for the best
finite-size collapse $m=F\!\left[(c-c^\ast)N^{1/\nu}\right]$ and compare its
residual against the best collapse obtainable after \emph{shuffling the size
labels}. A collapse is claimed only when it beats that null by a factor of two,
which prevents the two free parameters from manufacturing agreement.

\subsection{The consistency requirement}

Signatures are not allowed to accumulate merely by being individually
admissible. Each admissible signature contributes a boundary estimate, and the
estimates must agree to within $30\%$ of their mean before the family is
credited with more than one signature. Two diagnostics naming boundaries an
order of magnitude apart are evidence of an artefact, not of a transition; in
that case the signature count is capped at one. This rule was not decorative in
practice: it is what demoted our own flagship family from a transition claim to
a candidate boundary on a first pass, before the design flaw responsible was
found and corrected (\S\ref{sec:symmetry}).

\subsection{A scale path is not any second axis}

Finite-size analysis presumes that the second axis lists successive \emph{sizes
of one system}. Several natural experimental axes do not: a planted rank, a
number of experts, a generation count, and a label-noise level each change the
problem rather than enlarge it. Susceptibility and Binder curves can still be
computed along such an axis, and they usefully locate a threshold on the control
axis, but they cannot certify a thermodynamic statement. Families therefore
declare whether their second axis is a scale path, and those that do not are
capped at the outcome \emph{threshold located (no scale path)} regardless of how
clean their curves are. In this release that cap applies to
\AtlasCountThresholdLocatedNoScalePath{} families whose diagnostics would
otherwise have read as candidate boundaries.

\subsection{Sharpening and location are separate claims}

A susceptibility whose height diverges with size is evidence of criticality; it
does not say \emph{where} the boundary is. We therefore separate the two, and
the strongest verdict requires both. The resulting vocabulary is:

\begin{table}[tbp]
\centering\small
\begin{tabularx}{\linewidth}{@{}p{52mm}Y@{}}
\toprule
Outcome & Condition \\
\midrule
transition supported & $\ge2$ mutually consistent signatures \emph{and} a
resolved boundary location \\
critical response, boundary unresolved & diverging peak, but the location
extrapolation is inadmissible \\
candidate boundary & exactly one admissible signature \\
threshold located (no scale path) & the threshold is located on the control
axis, but the second axis is a problem parameter rather than a size, so no
finite-size statement is available \\
response shift (boundary off-grid) & the order parameter moves, and most sizes
are censored \\
response shift & the order parameter moves beyond its uncertainty; no
finite-size signature \\
no resolved response & the order parameter does not separate from its
uncertainty \\
\bottomrule
\end{tabularx}
\caption{Outcome vocabulary. Only the first row licenses phase language; the
remaining rows are reported with equal prominence rather than omitted.}
\label{tab:outcomes}
\end{table}

\subsection{Evidence grading}

Orthogonally to the outcome, each family carries a six-component evidence
vector --- semantics, replication, finite size, prediction, intervention, and
external validity --- each scored $0$, $1$ or $2$. Grade~A requires semantics
and replication both at $2$ with at least four components at $2$; grade~B
requires both at least $1$ with four components at least $1$; grade~C requires
replication at least $1$ with two components at least $1$. Replication reaches
$2$ only at twelve or more outer seeds. The present release runs eight outer
seeds per condition, which caps every family in this paper at grade~C by
construction: the finite-size evidence is real, and the replication tier
required for a stronger grade is deliberately not claimed.

\subsection{What is not tested here}

Two gates in the design remain open across the entire release and are stated
once rather than repeated per family. There is no frozen prospective re-run at
fresh seeds, so no result carries prospective confirmation. And no family is
evaluated on natural data or on a pretrained model, so external validity scores
zero throughout. Both are the natural next experiments rather than
qualifications buried in a limitations section.
